\section{Estimate Tab}
\label{sec:estimate}

The \app{Estimate} tab predicts beam time before a measurement. Given
an element, an electronic transition, an isotope chain, and beam
parameters, it computes per-peak signal rates, the on-peak time
required to reach a target statistical significance, the number of
voltage sweeps needed to cover the requested scan range, and a
grand-total run time. A \app{Spectrum only} switch suppresses the
timing pass and produces only the simulated lineshape --- useful when
planning the scan window itself.

The tab is a single surface with three resizable columns: the global
parameters on the left (Figure~\ref{fig:quickstart_estimate_tab}), the
isotope list in the middle, and a right column stacking the
\app{Run Options} group, the plot display, and the streaming log
(Figure~\ref{fig:quickstart_estimate_output}). The expand button on
the plot panel temporarily expands the plots over the whole tab;
click it again to bring the parameter columns back.

\subsection{Global parameters}

The left column holds settings shared by every isotope.

\begin{description}
    \item[\app{Element \& Transition}.]
        Element symbol, atomic number $Z$, the lower and upper levels
        of the electronic transition in cm$^{-1}$, and their total
        angular momenta $J_l$, $J_u$.
    \item[\app{Beam}.]
        Acceleration voltage (kV), ion charge state, and laser-ion
        geometry --- \app{collinear} or \app{anti-collinear}.
    \item[\app{Laser}.]
        Harmonic order and the setpoint of the \emph{fundamental}
        laser in cm$^{-1}$. As an alternative to a fixed setpoint,
        tick \app{Anchor to isotope}, choose an isotope and its state
        (g.s.\ or isomer), and enter the voltage offset \app{dV} from
        $V_0$ at which that resonance should sit; \denis{} then
        back-solves the laser frequency required to place the
        centroid at that voltage.
    \item[\app{Linewidth}.]
        Full-width-half-maximum used when drawing simulated spectra
        (MHz). Does not affect timing.
    \item[\app{Scan}.]
        Voltage range (full window, $\pm$ half on either side of
        $V_0$), step size, and dwell time per step (ms).
\end{description}

The checkable \app{Statistics (Timing Mode)} group is the master
switch for timing: unchecking it puts the whole tab in spectrum-only
mode (rates, backgrounds, and yields are hidden, and the scan step
size and dwell time are disabled). When checked it exposes:

\begin{itemize}
    \item \app{Required $\sigma$} --- target statistical significance per
        peak; defaults to $3\sigma$.
    \item \app{Uniform timing (global default)} --- when ticked, every
        measured peak gets the same on-peak time, computed for the
        strongest peak assuming full yield. Useful when a single dwell
        setting must cover the whole scan. Any isotope can override
        this locally with its own \app{Uniform timing} control.
\end{itemize}

\subsection{Isotope entries}

The middle column contains one expandable panel per isotope, added
with the \app{+ Add Isotope} button at the top. Panels are reordered
by dragging the hamburger handle and are color-striped by position so
the same isotope reads consistently across the tab. Each panel
exposes:

\begin{description}
    \item[Mass number $A$ and label] used in plot legends and log
        headers.
    \item[Nuclear spin $I$] integer or half-integer.
    \item[$\mu$ ($\mu_N$) and $Q$ (b)] magnetic dipole and electric
        quadrupole moments.
    \item[Isotope shift (MHz)] centroid offset from the reference
        isotope.
    \item[Yield (ions/s)] expected ion rate (a scientific-notation
        field).
    \item[Efficiency] one detected photon per
        $1/\eta$ ions; the log echoes this expanded as e.g.\
        ``$1$ photon per $500$ ions''.
    \item[Peaks to measure] how many of the strongest hyperfine peaks
        enter the timing estimate (default 1; raise it when several
        components are diagnostic).
    \item[Background] either a flat rate (Hz), or a continuous rate
        combined with a TOF-gate width (\si{\micro\second}) and the
        dwell time, which then collapses to an effective rate
        automatically.
\end{description}

Exactly one isotope is marked as the \app{Reference}. Its hyperfine
$A$ and $B$ constants (entered in MHz, lower- and upper-state
separately) are used to scale the constants of every other isotope
via the Schmidt-moment ratio. Tick \app{Override A/B constants} on a
non-reference isotope to bypass scaling and supply explicit values ---
useful when high-quality measured constants are already available
(this is the case shown for \textsuperscript{51}V in
Figure~\ref{fig:quickstart_estimate_tab}).

\paragraph{Schmidt-moment auto-calculation.}
For an isotope whose $\mu$ is unknown, tick \app{Use Schmidt moment}
to derive it from the shell model. The single-particle case takes a
nucleon type (\code{p}/\code{n}) and an orbital (e.g.\ \code{f7/2}).
For odd-odd nuclei, choose the two-particle variant and supply two
orbital configurations together with the resulting spin $I$. The
unquenched Schmidt $\mu$ is auto-filled into the panel's $\mu$ field
(undoable with \code{Ctrl+Z}); the configured $g_s$ quenching factor
is applied during the run, and the log prints the configuration
string next to the value actually used. The same calculator is
available standalone under \menu{Tools $\rightarrow$ Schmidt Moment
Calculator} (Section~\ref{sec:tools}).

\paragraph{Isomers.}
Tick \app{Has isomer} to expand a sub-panel for an isomeric state.
The isomer carries its own spin, moments, isotope shift, and one of
three yield specifications: an explicit \app{Yield} (independent of
the ground state), a \app{Production ratio} in $[0,1]$ that splits
the parent's yield, or a \app{Spin distribution} $\sigma$ that
estimates the split from spin statistics --- the latter is useful for
surface-ion sources where high-spin states are populated
preferentially. The isomer contributes its own peaks to the simulated
spectrum and its own row to the timing table.

\subsection{Running the estimation}

Click \app{Run Estimation} in the \app{Run Options} group (right
column). The log streams in real time and, on completion, contains
in order:

\begin{enumerate}
    \item A header summarising the inputs and, in anchor mode, the
        derived laser setpoint.
    \item Per isotope: signal rate, on-peak time per measured peak,
        the number of sweeps required to accumulate enough statistics
        on each peak, and a subtotal in hours and in 8\,h shifts.
    \item A \app{Nuclear parameters} summary table with the $\mu$
        values actually used (config or Schmidt) and the resulting
        $A$, $B$ constants.
    \item The grand-total run time across all isotopes.
    \item Two peak lists --- one sorted by intensity, one by voltage
        --- covering every hyperfine component of every isotope and
        isomer.
\end{enumerate}

The log is also written verbatim to a \code{.log} file in the run's
output folder (\code{cls\_<timestamp>} under the chosen
\app{Output directory}); the YAML used for the run is preserved
alongside it for reproducibility. The run additionally persists its
plot arrays and peak table (\code{estimate\_results.npz},
\code{peaks.csv}), and the save file records the run folder --- so a
saved session restores its plots, peak table, and log on load without
re-running.

\subsection{Simulated spectra}

The plot panel shows a simulated voltage spectrum per isotope on the
configured scan range, with vertical dashed lines at each hyperfine
peak position (Figure~\ref{fig:quickstart_estimate_output}b). The
\app{View} drop-down switches between \app{Individual Spectra} (one
panel per isotope), \app{Combined Overview} (all isotopes overlaid),
and \app{Peak List} (a sortable table of every hyperfine component).
\app{Edit Plot} opens the matplotlib style editor; \app{Open PDF}
reveals the exported PDF in the system file viewer.

\subsection{Loading a previous estimation}

\app{Load Estimation\dots} (next to \app{Run Estimation}) reopens any
previous \code{cls\_<timestamp>} run folder for viewing --- plots,
peak table, and log are restored without re-running and without
creating new output. Runs from before the \code{.npz} persistence
existed hold only PDFs; their pages are rendered onto the live
matplotlib canvas, so pan/zoom/save still work.
